\section{Operational semantics}
\label{sec:opsem}

\subsection{Values and environments}

Values, environments, and evaluation results are given in \figref{values}. For a statement $s$, the result
$\assigns{\rho}$ represents \emph{normal completion} of $s$, with accumulated environment $\rho$; the result
$\returns{v}$ represents \emph{abrupt completion} of $s$, with return value $v$; \Replaced{and the result
$\aborts{\kappa}$ represents abrupt completion of the whole run, with termination kind $\kappa$}{and
$\aborts{\kappa}$ abrupt completion with an error}.
Composing statements sequentially composes their results according to the following partial monoid, analogous
to sequential composition of result types:
{\small
\begin{align*}
\assigns{\rho} \override \assigns{\rho'} & = \assigns{(\rho \override \rho')} \\
\assigns{\rho} \override \returns{v} & = \returns{v} \\
\assigns{\rho} \override \aborts{\kappa} & = \aborts{\kappa}
\end{align*}}

with $\assigns{\varnothing}$ as left and right unit.

\begin{figure}
\small
\begin{center}
\begin{tabular}{rcll}
\grammarGroup{Values} \\
$v$ & $::=$ & $\ell$ & (Literal) \\
    & $\mid$ & $\exList{v_1, \ldots, v_n}$ & (List) \\
    & $\mid$ & $\exTuple{v_1, \ldots, v_n}$ & (Tuple) \\
    & $\mid$ & $\exDict{\delta}$ & (Dict) \\
    & $\mid$ & $\exLamClosure[\Gamma]{\rho}{\seq{x}}{e}$ & (LambdaClosure) \\
    & $\mid$ & $\exDefClosure[\Gamma]{\rho}{\seq{d}}{i}$ & (DefClosure) \\
    & $\mid$ & $\exObj{\Gamma}{C}{\rho}$ & (Object) \\
    & $\mid$ & $\modStub{q}$ & (ModuleStub) \\
    & $\mid$ & $\modLoaded{q}{\rho}$ & (LoadedModule) \\
    & $\mid$ & $\clsEntry{\Gamma}{q.c}{\seq{x}}{c'}$ & (Class) \\
    & $\mid$ & $\primVal{\primFn}$ & (Primitive)
   \\[3mm]
\grammarGroup{Environments} \\
$\rho$ & $::=$ & $\set{\bind{x_1}{v_1}, \ldots, \bind{x_n}{v_n}}$ & 
   \\[3mm]
\grammarGroup{Dictionary entries} \\
$\delta$ & $::=$ & $\bind{w_1}{v_1} \cons \ldots \cons \bind{w_n}{v_n}$ &
   \\[3mm]
\grammarGroup{Results} \\
$r$ & $::=$ & $\assigns{\rho}$ & (Assigns) \\
  & $\mid$ & $\returns{v}$ & (Returns) \\
  & $\mid$ & $\aborts{\kappa}$ & (Aborts)
   \\[3mm]
\grammarGroup{Outcomes} \\
$o$ & $::=$ & $\val{v}$ & (Val) \\
  & $\mid$ & $\aborts{\kappa}$ & (Aborts)
   \\[3mm]
\grammarGroup{Match results} \\
$m$ & $::=$ & $\matchOk{\rho}$ & (Match) \\
  & $\mid$ & $\noMatch$ & (NoMatch)
   \\[3mm]
\grammarGroup{Termination kinds} \\
$\kappa$ & $::=$ & $\errType \mid \errIndex \mid \errKey \mid \errZeroDiv$ &  \\
    & $\mid$ & $\errAssertNoMsg \mid \errAssert{w} \mid \errAttr \mid \errExit{n}$ &
\end{tabular}
\end{center}
\caption{Values, environments and evaluation results}
\label{fig:values}
\end{figure}


\Added{The primitive $\primVal{\primFn}$ is how $\mathcal{P}$ supplies a member that can be called, since a
supplied member has no body to evaluate and $\primFn$ gives its meaning directly. No rule applies a
primitive, and \ruleName{eval-call-nonfun} does not cover one, so calling a primitive has no derivation
(\issue{132}, \issue{47}).}

A value is well-formed at a context entry (\figref{well-formed-env}): ordinary values at
$\statusDefAssigned$, module references and class entries at the corresponding entry. An environment is
well-formed at a context that binds its variables at corresponding entries.

\begin{figure}
\flushleft \shadebox{$\vdash v$}
\begin{mathpar}
\small
\inferrule*[lab={\ruleName{val-lit}}]
{
  \strut
}
{
  \vdash \ell
}
\and
\inferrule*[lab={\ruleName{val-list}}]
{
  \vdash v_i \quad (\forall i)
}
{
  \vdash \exList{v_1, \ldots, v_n}
}
\and
\inferrule*[lab={\ruleName{val-tuple}}]
{
  \vdash v_i \quad (\forall i)
}
{
  \vdash \exTuple{v_1, \ldots, v_n}
}
\and
\inferrule*[lab={\ruleName{val-lambda}}]
{
  \Gamma \vdash \rho \\
  \Gamma \runion \statusDefAssigned(\seq{x}) \vdash e
}
{
  \vdash \exClosure[\Gamma]{\rho}{\seq{x}}{e}
}
\and
\inferrule*[lab={\ruleName{val-def}}]
{
  \Gamma \vdash \rho \\
  d_j = \exDef{f_j}{\vec{x}_j}{b_j} \quad (\forall j) \\
  \Gamma \runion \statusDefAssigned(\seq{f}) \runion \statusDefAssigned(\seq{x}_j) \runion
  \statusNotDefAssigned(\assignsS(b_j) \setminus \seq{x}_j) \vdash b_j : R_j \quad (\forall j) \\
  \seq{f} \text{ distinct}
}
{
  \vdash \exDefClosure[\Gamma]{\rho}{\seq{d}}{i}
}
\and
\inferrule*[lab={\ruleName{val-obj}}]
{
  \Gamma \vdash \rho
}
{
  \vdash \exObj[\Gamma]{d}{\rho}
}
\end{mathpar}

\vspace{2mm}

\flushleft \shadebox{$\Gamma \vdash \rho$}
\begin{mathpar}
\small
\inferrule*[lab={\ruleName{env-var}}]
{
  x : a \in \Gamma \\
  \vdash v
}
{
  \Gamma \vdash \set{x \mapsto v}
}
\and
\inferrule*[lab={\ruleName{env-module}}]
{
  x : \modEntry{q} \in \Gamma \\
  \Delta \vdash \rho'
}
{
  \Gamma \vdash \set{x \mapsto \exObj[\Delta]{\bot}{\rho'}}
}
\and
\inferrule*[lab={\ruleName{env-class}}]
{
  x : \clsEntry{q}{\vec{y}}{c} \in \Gamma
}
{
  \Gamma \vdash \set{x \mapsto \clsEntry{q}{\vec{y}}{c}}
}
\and
\inferrule*[lab={\ruleName{env-cons}}]
{
  \Gamma \vdash \set{x \mapsto w} \\
  \Gamma \vdash \rho
}
{
  \Gamma \vdash \set{x \mapsto w} \runion \rho
}
\end{mathpar}
\caption{Well-formed values and environments}
\label{fig:well-formed-env}
\end{figure}


The rules are given in a monadic metalanguage over outcomes (\figref{values}), in which
$\mlet{v}{o}{\ldots}$ binds the value of an outcome; for example, \ruleName{eval-binop} concludes
$\mlet{v'}{o}{\hat{\odot}(v, v')}$, which is $\hat{\odot}(v, v')$ where $o$ is $\val{v'}$, and
$\aborts{\kappa}$ where $o$ is $\aborts{\kappa}$. A final
premise needs no propagation rule; earlier premises have one each, recording which subexpressions succeeded
before the failing one was reached (\figref{eval-fails}). The statement, sequence and module judgements use
outcomes of the same form, carrying a result, a sequence of values or an environment in place of a value.

\subsection{Auxiliary definitions}

The strict operators are given by metafunctions $\hat{\odot}$ and $\hat{\ominus}$ (\figref{operators}), in
terms of structural equality $\eqOp$ and membership $\memOp$ (\figrefTwo{equality}{membership}). They have
the meanings they have in Python and yield $\aborts{\kappa}$ where Python raises. All three are undefined
where no case applies, and where one of them is undefined the expression has no derivation, which is where
\PurePy declines to give the operation the meaning Python gives it; a non-Boolean test is the same case
(\secref{no-coercions}, \issue{162}).

\begin{figure}
\small
\begin{center}
\begin{tabular}{>{\raggedright\arraybackslash}p{0.22\textwidth}>{\raggedright\arraybackslash}p{0.62\textwidth}}
\hline
$\odot$ & $\hat{\odot}(v, v')$ \\
\hline
\pyinline{==} & $\val{\pyinline{False}}$ if both operands are \pyinline{nan}, and $\val{\eqOp(v, v')}$ otherwise \\
\pyinline{!=} & $\val{\pyinline{True}}$ if both operands are \pyinline{nan} or $\eqOp(v, v') = \pyinline{False}$, and $\val{\pyinline{False}}$ if $\eqOp(v, v') = \pyinline{True}$ \\
\pyinline{in} & $\val{\memOp(v', v)}$ \\
\pyinline{not in} & $\val{\pyinline{True}}$ if $\memOp(v', v) = \pyinline{False}$, and $\val{\pyinline{False}}$ if $\memOp(v', v) = \pyinline{True}$ \\
\pyinline{<}\quad\pyinline{<=}\quad\pyinline{>}\quad\pyinline{>=} & \todo{Ordering, as in Python} \\
\pyinline{+}\quad\pyinline{-}\quad\pyinline{*}\quad\pyinline{/}\quad\pyinline{//}\quad\pyinline{\%}\quad\pyinline{**} & \todo{Arithmetic, as in Python. Float precision is left to the implementation. $\aborts{\errZeroDiv}$ where the divisor is zero, and $\aborts{\errType}$ where an operand is not a number} \\
\hline
\end{tabular}


\vspace{2mm}

\input{functions/unop}
\end{center}
\caption{Operator meanings. An operator returns an outcome}
\label{fig:operators}
\end{figure}

\begin{figure}
\small
\begin{align*}
\eqOp(\ell, \ell') &{}= \pyinline{True} && \text{if } \ell = \ell' \\
\eqOp(\ell, \ell') &{}= \pyinline{False} && \text{if } \ell \neq \ell' \text{ and } \ell, \ell' \text{ have the same kind} \\
\\
\eqOp(\exList{v_1, \ldots, v_n}, \exList{v_1', \ldots, v_n'}) &{}= \textstyle\bigwedge_{1 \leq i \leq n} \eqOp(v_i, v_i') \\
\eqOp(\exTuple{v_1, \ldots, v_n}, \exTuple{v_1', \ldots, v_n'}) &{}= \textstyle\bigwedge_{1 \leq i \leq n} \eqOp(v_i, v_i') \\
\eqOp(\exList{\seq{v}}, \exList{\seq{v}\,'}) = \eqOp(\exTuple{\seq{v}}, \exTuple{\seq{v}\,'}) &{}= \pyinline{False} && \text{if } \seq{v}, \seq{v}\,' \text{ differ in length} \\
\\
\eqOp(d, d') &{}= \textstyle\bigwedge_{w \in \dom(d)} \eqOp(d(w), d'(w)) && \text{if } \dom(d) = \dom(d') \\
\eqOp(d, d') &{}= \pyinline{False} && \text{if } \dom(d) \neq \dom(d') \\
\\
\eqOp(\exObj{\Gamma}{C}{\rho}, \exObj{\Gamma'}{C}{\rho'}) &{}= \textstyle\bigwedge_{x \in \fields(C)} \eqOp(\rho(x), \rho'(x)) \\
\eqOp(\exObj{\Gamma}{C}{\rho}, \exObj{\Gamma'}{C'}{\rho'}) &{}= \pyinline{False} && \text{if } C \neq C'
\end{align*}
\caption{Structural equality; undefined on operands of differing kinds and on closures}
\label{fig:equality}
\end{figure}

\begin{figure}
\small
\begin{align*}
\memOp(\exList{v_1, \ldots, v_n}, v) &{}= \memElems(v_1 \cons \ldots \cons v_n, v) \\*
\memOp(\exTuple{v_1, \ldots, v_n}, v) &{}= \memElems(v_1 \cons \ldots \cons v_n, v) \\*
\\
\memOp(\delta, w) &{}= \pyinline{True} && \text{if } w \in \dom(\delta) \\*
\memOp(\delta, w) &{}= \pyinline{False} && \text{if } w \notin \dom(\delta) \\*
\\
\memOp(w', w) &{}= \pyinline{True} && \text{if } w' = w_1 \concat w \concat w_2 \text{ for some } w_1, w_2 \\*
\memOp(w', w) &{}= \pyinline{False} && (\text{otherwise})
 \\
\\
\memElems(\varepsilon, v) &{}= \pyinline{False} \\*
\memElems(v' \cons \seq{v}, v) &{}= \pyinline{True} && \text{if } \eqOp(v, v') = \pyinline{True} \\*
\memElems(v' \cons \seq{v}, v) &{}= \memElems(\seq{v}, v) && \text{if } \eqOp(v, v') = \pyinline{False}

\end{align*}
\caption{Membership, undefined where no case applies}
\label{fig:membership}
\end{figure}


\Added{The remaining metafunctions the rules use are collected in \figref{eval-aux}: $\outcome$ reads a
block's result as the outcome of calling it, $\iterOp$ gives the sequence a generator draws from, $\getitem$
interprets subscripting, $\dispatch$ selects the case a $\pyinline{match}$ takes, and $\entries$ and
$\update$ build a dictionary from its entries.}

\begin{figure}
\small
\begin{align*}
\elems(\exList{v_1, \ldots, v_n}) &{}= v_1 \ldots v_n \\
\elems(\exTuple{v_1, \ldots, v_n}) &{}= v_1 \ldots v_n \\
\\
\added{\iterOp(v)} &{}\added{{}= \elems(v)} \\
\\
\added{\getitem(v, n)} &{}\added{{}= \val{v_{n+1}}} && \added{\text{if } \elems(v) = v_1 \ldots v_m \text{ and } 0 \leq n < m} \\
\added{\getitem(v, n)} &{}\added{{}= \val{v_{m+n+1}}} && \added{\text{if } \elems(v) = v_1 \ldots v_m \text{ and } {-m} \leq n < 0} \\
\added{\getitem(v, n)} &{}\added{{}= \fails{\errIndex}} && \added{\text{if } \elems(v) = v_1 \ldots v_m \text{ and } (n < -m \text{ or } n \geq m)} \\
\added{\getitem(v, v')} &{}\added{{}= \fails{\errType}} && \added{(\text{otherwise})} \\
\\
\outcome(\returns{v}) &= \val{v} \\
\outcome(\assigns{\rho}) &= \val{\pyinline{None}} \\
\outcome(\fails{\kappa}) &= \fails{\kappa} \\
\\
\env(\rho, x) &= \rho' && \text{where } \rho(x) = \modLoaded{q'}{\rho'} \\
\env(\rho, q.x) &= \rho' && \text{where } \env(\rho, q)(x) = \modLoaded{q'}{\rho'} \\
\\
\resolveClass(\rho, c) &= \rho(c) \\
\resolveClass(\rho, q.c) &= \env(\rho, q)(c) \\
\\
\dispatch(\rho, v, \varepsilon, \varepsilon) &= (\varnothing, \exPass) \\
\dispatch(\rho, v, p \cons \seq{p}, b \cons \seq{b}) &= (\rho', b) && \text{if } \rho, p, v \matches \rho' \\
\dispatch(\rho, v, p \cons \seq{p}, b \cons \seq{b}) &= \dispatch(\rho, v, \seq{p}, \seq{b}) && \text{if } \rho, p, v \notmatches \\
\\
\entries(\delta, \varepsilon) &{}= \delta \\
\entries(\delta, \bind{w}{v} \cons \delta') &{}= \entries(\update(\delta, w, v), \delta') \\
\\
\update(\varepsilon, w, v) &{}= \bind{w}{v} \\
\update(\bind{w}{u} \cons \delta, w, v) &{}= \bind{w}{v} \cons \delta \\
\update(\bind{w'}{u} \cons \delta, w, v) &{}= \bind{w'}{u} \cons \update(\delta, w, v) && \text{if } w \neq w'
\end{align*}
\caption{Metafunctions used by block and expression evaluation}
\label{fig:eval-aux}
\end{figure}


\subsection{Pattern matching}

The pattern-matching judgement $\rho, p, v \matches m$ (\figref{match}) is a deterministic partial
function. \Added{A \emph{match result} $m$ is either $\matchOk{\rho'}$, a match binding $\rho'$, or
$\noMatch$, a failure to match, on which $\dispatch$ moves to the next case. Having no derivation is a third
outcome, distinct from both: the match has no meaning, and \PurePy gives the run no result. That is where a
pattern and a value are not comparable, so \PurePy can say neither that they match nor that they do not; a
list pattern against a tuple is such a case (\issue{151}). Sub-results compose with $\override$, matching
sequentially:}
{\small
\begin{align*}
\added{\matchOk{\rho} \override \matchOk{\rho'}} & \added{{}= \matchOk{(\rho \disjunion \rho')}} \\
\added{m \override \noMatch = \noMatch \override m} & \added{{}= \noMatch}
\end{align*}}

\Added{with $\matchOk{\varnothing}$ as left and right unit.}
A constructor pattern matches any object whose class has the
pattern's class among its ancestors; sub-patterns bind the fields of the pattern's class. An object records
its class entry, whose qualified name is the identity of the class: since a module declares each class name at
most once, the entry comparison in \ruleName{eval-pat-constr} amounts to a comparison of qualified names.

\begin{figure}
\flushleft \shadebox{$\rho \vdash p, v \matches \rho'$}
\begin{mathpar}
\small
\inferrule*[lab={\ruleName{pat-lit}}]
{
  \strut
}
{
  \rho \vdash \ell, \ell \matches \varnothing
}
\and
\inferrule*[lab={\ruleName{pat-neg-lit}}]
{
  \strut
}
{
  \rho \vdash -n, -n \matches \varnothing
}
\and
\inferrule*[lab={\ruleName{pat-var}}]
{
  \strut
}
{
  \rho \vdash x, v \matches \set{x \mapsto v}
}
\and
\inferrule*[lab={\ruleName{pat-wild}}]
{
  \strut
}
{
  \rho \vdash \patWild, v \matches \varnothing
}
\and
\inferrule*[lab={\ruleName{pat-list}}]
{
  \rho \vdash p_i, v_i \matches \rho_i \quad (\forall i)
}
{
  \rho \vdash \patList{p_1, \ldots, p_n}, \exList{v_1, \ldots, v_n} \matches \rho_1 \disjunion \ldots \disjunion \rho_n
}
\and
\inferrule*[lab={\ruleName{pat-tuple}}]
{
  \rho \vdash p_i, v_i \matches \rho_i \quad (\forall i)
}
{
  \rho \vdash \patTuple{p_1, \ldots, p_n}, \exTuple{v_1, \ldots, v_n} \matches \rho_1 \disjunion \ldots \disjunion \rho_n
}
\and
\inferrule*[lab={\ruleName{pat-as}}]
{
  \rho \vdash p, v \matches \rho'
}
{
  \rho \vdash p\;\pyinline{as}\;x, v \matches \rho' \disjunion \set{x \mapsto v}
}
\and
\inferrule*[lab={\ruleName{pat-class}}]
{
  \fields(\Gamma_q(c)) = x_1, \ldots, x_n, x_{n+1}, \ldots, x_{n+m} \\
  \rho \vdash p_i, \rho_o(x_i) \matches \rho_i \quad (\forall i.\, 1 \le i \le n) \\
  \rho \vdash p_{n+j}, \rho_o(y_j) \matches \rho_{n+j} \quad (\forall j.\, 1 \le j \le m)
}
{
  \rho \vdash \patClass{c}{p_1, \ldots, p_n, y_1 \mathord{=} p_{n+1}, \ldots, y_m \mathord{=} p_{n+m}}, \exObj{q.c}{\rho_o} \matches \rho_1 \disjunion \ldots \disjunion \rho_{n+m}
}
\end{mathpar}
\caption{Pattern matching}
\label{fig:match}
\end{figure}


\subsection{Blocks and expressions}

Block evaluation $\rho, b \Rightarrow r$\Added{, for any $\rho : \Gamma$ and any $b$ well-formed in $\Gamma$,}
is given in \figref{eval-block-outcome}. Function definitions delimit
the scope of sequential-composition effects, so as black boxes functions are always pure. The metafunction $\outcome$ (\figref{eval-aux}) reads a block's result as
the outcome of calling it, $\val{\pyinline{None}}$ where the block falls off the end.

Expression evaluation $\rho, e \Rightarrow o$ (\figrefTwo{eval-outcome}{eval-outcome-ctd})\Added{, for any $\rho : \Gamma$ and any $e$ well-formed in $\Gamma$,} yields an
outcome, either $\val{v}$ or $\aborts{\kappa}$ (\figref{values}). \Replaced{The termination kinds other than $\errExit{n}$ are named after the
exceptions the same program raises under Python in the same circumstances.}{The error kinds are designed to
correspond to the error the same program would raise under Python in the same circumstances.} The operators
\pyinline{and} and \pyinline{or} evaluate their second operand only when the first leaves the result open, and
a conditional expression requires its condition to be \pyinline{True} or \pyinline{False}, as
\ruleName{eval-if} does. A comprehension evaluates its body once for each binding its qualifiers produce, and
concatenates the results; a generator draws from $\iterOp$ (\figref{eval-aux}), which is
\Replaced{defined on lists, tuples, strings and dictionaries, a dictionary yielding its keys as in Python.
These are the only values a generator can draw from.}{defined on lists and tuples only, since \PurePy has no
wider notion of iteration.}

\begin{figure}
\small
\begin{mathpar}
\inferrule*[lab={\ruleName{eval-cond-true}}]
{
  \rho, e \Rightarrow \val{\pyinline{True}} \\
  \rho, e_1 \Rightarrow o
}
{
  \rho, \exCondExpr{e_1}{e}{e_2} \Rightarrow o
}
\end{mathpar}
\caption{\todo{Expression $e$ evaluates to outcome $o$ under environment $\rho : \Gamma$, where $o$ is
$\val{v}$ or $\fails$ (\issue{63})}}
\label{fig:eval-outcome}
\end{figure}


\subsection{Modules and programs}

Module and program evaluation are given in \figref{eval-modules}; import evaluation in \figref{eval-import-prefix}; importing names from a loaded module in
\figref{eval-imports}. Extend (Definition~\ref{def:extend}) extends to module values by reading its equations with a loaded reference $\modLoaded{q}{\rho}$ carrying an environment in place of a context. The loading judgement $q \Rightarrow_{\mathsf{load}} o$ evaluates the import prefix
of $q$ and then its statements, mirroring the well-formedness judgement $\vdash_{\mathsf{load}} q : \Gamma$
(\figref{well-formed-module}); as with the static rules, the mutual
recursion between loading and evaluation is grounded by program well-formedness, which rules out
import cycles (\secref{modules-programs}).
Loading takes a single module; the ancestors of an imported module are loaded by the
loads-as judgement for an import, and directly as premises of \ruleName{eval-from-import} for a
from-import.

\Deleted{The correspondence is exact: if $\vdash_{\mathsf{load}} q : \Gamma$ and $q \Rightarrow_{\mathsf{load}} \val{\rho}$ then $\vdash \rho : \Gamma$ (\figref{well-formed-env}).} \begin{theorem}[\Added{Load preservation}]
\Added{If $\vdash_{\mathsf{load}} q : \Gamma$ and $q \Rightarrow_{\mathsf{load}} \val{\rho}$ then $\vdash \rho : \Gamma$.}
\end{theorem}

\noindent Since
imports precede statements, every module a statement can reference is loaded before the statement runs, and
\Replaced{loading is deterministic, so}{since loading has no observable effects, reloading a module is
unobservable; the semantics therefore needs no cache of loaded modules.} \Added{the semantics needs no cache
of loaded modules. The rule \ruleName{eval-module} puts $\mathcal{P}(\pyinline{builtins})$ under a module
body's environment, mirroring \ruleName{module} (\secref{wf-modules}).}\Deleted{These rules use the environment $\mathcal{P}(\pyinline{builtins})$ binding the names of
\figref{predefined}.}
\Replaced{A program evaluates to a \emph{termination kind} (\figref{values}): $\errExit{0}$ if the main module
loads, and otherwise the kind it aborted with. $\errExit{n}$ is a run that exited with status $n$; the remaining kinds are named after the exceptions the
same program raises under Python. An implementation must agree on which kind a run yields, though how it
reports one is not prescribed.}{A program evaluates to
its exit status: a run whose main module loads successfully yields $0$, and one whose main module fails yields
$1$, as Python does for an uncaught exception.}

\begin{figure}
\begin{subfigure}{\textwidth}
\flushleft \shadebox{$\mu, \rho, m \Rightarrow \mu', \rho'$}
\begin{mathpar}
\small
\inferrule*[lab={\ruleName{eval-nil}}]
{
  \strut
}
{
  \mu, \rho, \varepsilon \Rightarrow \mu, \varnothing
}
\and
\inferrule*[lab={\ruleName{eval-import}}]
{
  \mu, q \Rightarrow_{\mathsf{load}} \mu' \\
  q = x_1.x_2.\ldots.x_{n+1} \\
  \mu', \rho \runion \set{\bind{x_1}{\exMod{\Gamma_{x_1}}{\mu'(x_1)}}}, m \Rightarrow \mu'', \rho'
}
{
  \mu, \rho, \exImport{q} \cons m \Rightarrow \mu'', \rho'
}
\and
\inferrule*[lab={\ruleName{eval-from-import}}]
{
  \mu, q \Rightarrow_{\mathsf{load}} \mu' \\
  \mu', q, \vec{x} \importsR \mu'', \rho' \\
  \mu'', \rho \runion \rho', m \Rightarrow \mu''', \rho''
}
{
  \mu, \rho, \exFromImport{q}{\vec{x}} \cons m \Rightarrow \mu''', \rho''
}
\and
\inferrule*[lab={\ruleName{eval-statement}}]
{
  \rho, s \Rightarrow \assigns{\rho'} \\
  \mu, \rho \runion \rho', m \Rightarrow \mu', \rho''
}
{
  \mu, \rho, s \cons m \Rightarrow \mu', \rho' \runion \rho''
}
\and
\inferrule*[lab={\ruleName{eval-class}}]
{
  C = \clsEntry{\Gamma}{q}{\ownFields(\phi)}{\bot} \\
  \mu, \rho \runion \set{\bind{c}{C}}, m \Rightarrow \mu', \rho'
}
{
  \mu, \rho, \exDataclass{c}{\phi} \cons m \Rightarrow \mu', \set{\bind{c}{C}} \runion \rho'
}
\and
\inferrule*[lab={\ruleName{eval-class-extend}}]
{
  C = \clsEntry{\Gamma}{q}{\ownFields(\phi)}{c'} \\
  \mu, \rho \runion \set{\bind{c}{C}}, m \Rightarrow \mu', \rho'
}
{
  \mu, \rho, \exDataclassExtend{c}{c'}{\phi} \cons m \Rightarrow \mu', \set{\bind{c}{C}} \runion \rho'
}
\end{mathpar}
\subcaption{Module body $m$ of module $q$ evaluates under store $\mu$ and environment $\rho : \Gamma$}
\end{subfigure}

\vspace{2mm}

\begin{subfigure}{\textwidth}
\flushleft \shadebox{$\mathcal{M} \Rightarrow n$}
\begin{mathpar}
\small
\inferrule*[lab={\ruleName{eval-program}}]
{
  \varnothing, \varnothing, (\exAssign{\pymath{\_\_name\_\_}}{\pymath{"\_\_main\_\_"}}) \cons \mathcal{M}(\pymath{\_\_main\_\_}) \Rightarrow \mu, \rho
}
{
  \mathcal{M} \Rightarrow 0
}
\end{mathpar}
\subcaption{Program $\mathcal{M}$ evaluates to exit code $n$}
\end{subfigure}
\caption{Module and program evaluation}
\label{fig:eval-modules}
\end{figure}

\begin{figure}
\begin{subfigure}{\textwidth}
\flushleft \shadebox{$\mu, q \Rightarrow_{\mathsf{load}} \mu'$}
\begin{mathpar}
\small
\inferrule*[lab={\ruleName{load-cached}}]
{
  q \in \dom(\mu)
}
{
  \mu, q \Rightarrow_{\mathsf{load}} \mu
}
\and
\inferrule*[lab={\ruleName{load-unqualified}}]
{
  x \notin \dom(\mu) \\
  \mu, \varnothing, (\exAssign{\pymath{\_\_name\_\_}}{\nameOf(x)}) \cons \mathcal{M}(x) \Rightarrow \mu', \rho
}
{
  \mu, x \Rightarrow_{\mathsf{load}} \mu' \runion \set{\bind{x}{\rho}}
}
\and
\inferrule*[lab={\ruleName{load-qualified}}]
{
  q.x \notin \dom(\mu) \\
  \mu, q \Rightarrow_{\mathsf{load}} \mu' \\
  \mu', \varnothing, (\exAssign{\pymath{\_\_name\_\_}}{\nameOf(q.x)}) \cons \mathcal{M}(q.x) \Rightarrow \mu'', \rho
}
{
  \mu, q.x \Rightarrow_{\mathsf{load}} \mu'' \runion \set{\bind{q.x}{\rho}}
}
\end{mathpar}
\subcaption{Module $q$ is loaded into the store $\mu$}
\end{subfigure}

\vspace{2mm}

\begin{subfigure}{\textwidth}
\flushleft \shadebox{$\mu, q, \vec{x} \importsR \mu', \rho$}
\begin{mathpar}
\small
\inferrule*[lab={\ruleName{imp-nil}}]
{
  \strut
}
{
  \mu, q, \varepsilon \importsR \mu, \varnothing
}
\and
\inferrule*[lab={\ruleName{imp-var}}]
{
  x \in \dom(\mu(q)) \\
  \mu, q, \vec{x} \importsR \mu', \rho
}
{
  \mu, q, x \cons \vec{x} \importsR \mu', \set{\bind{x}{\mu(q)(x)}} \runion \rho
}
\and
\inferrule*[lab={\ruleName{imp-class}}]
{
  x \notin \dom(\mu(q)) \\
  q.x \notin \dom(\mathcal{M}) \\
  \mu, q, \vec{x} \importsR \mu', \rho
}
{
  \mu, q, x \cons \vec{x} \importsR \mu', \rho
}
\and
\inferrule*[lab={\ruleName{imp-module}}]
{
  x \notin \dom(\mu(q)) \\
  q.x \in \dom(\mathcal{M}) \\
  \mu, q.x \Rightarrow_{\mathsf{load}} \mu' \\
  \mu', q, \vec{x} \importsR \mu'', \rho
}
{
  \mu, q, x \cons \vec{x} \importsR \mu'', \set{\bind{x}{\mu'(q.x)}} \runion \rho
}
\end{mathpar}
\subcaption{Names $\vec{x}$ imported from module $q$ produce bindings $\rho$}
\end{subfigure}
\caption{Module loading and importing}
\label{fig:eval-load}
\end{figure}

